% !TEX root = ../main.tex
\chapter{Forwards and Futures}\label{ch:futures}
\begin{paracol}{2}

\begin{sources}
\begin{tabularx}{\linewidth}{@{}l l L@{}}
\toprule
\textbf{Segment} & \textbf{Length} & \textbf{Content}\\
\midrule
\seg{18}{1}{L18.1 Preview: FT: Argentina in court to fight debt ruling} & 4:56 & Holdouts, pari passu and collective action clauses.\\
\seg{18}{2}{L18.2 Banking as advance clearing} & 5:55 & Changed expectations cause cash flows today: the survival constraint again.\\
\seg{18}{3}{L18.3 Forwards versus futures} & 13:41 & Forward rate agreements as matched book; speculators and futures; $f>$ futures $>\E[r]$.\\
\seg{18}{4}{L18.4 Futures contracts, fluctuations II} & 14:07 & Convergence; the value of a forward contract over its life.\\
\seg{18}{5}{L18.5 Futures contracts, fluctuations II} & 6:07 & Futures marked to market daily: cash flows all along.\\
\seg{18}{6}{L18.6 Cash and carry arbitrage, defined} & 7:28 & Stigum's example; full carry pricing; implied repo rate.\\
\seg{18}{7}{L18.7 Cash and carry arbitrage, explained as liquidity risk} & 5:26 & Fully hedged, but not in cash flow.\\
\seg{18}{8}{L18.8 Cash and carry arbitrage as counterparty risk} & 2:49 & An alternative explanation.\\
\seg{18}{9}{L18.9 Cash and carry arbitrage, as natural banking business} & 3:22 & Banks' comparative advantage in liquidity risk.\\
\bottomrule
\end{tabularx}

\smallskip
\textbf{Files.} Transcripts \file{materials/transcripts/L18-P*.txt}. (The two segments 4 and 5 carry the same title on the playlist, but their content differs.) Mehrling's notes \mn{18}{1}.
\textbf{Reading.} M.~Stigum, \emph{The Money Market}, pp.~718--722 (cash and carry), as cited in class; J.~R.~Hicks, \emph{Value and Capital} (1939), ch.~XI; FT press cutting on Argentina (all in \file{materials/readings/}).
\end{sources}

\section*{Motivation}
The last section of the course is called ``banking as advance clearing'' \ts{18}{2}{0:03}. The first section showed the payment system: cash inflows and
outflows balanced each day, deficits pushed into the future, stress showing up as an interest rate, and defaults when it gets big enough \ts{18}{2}{0:26}.
The economy is coherent not because it is an intertemporal general equilibrium but because the clearing constraint forces everyone to pay attention to their place in
the system \ts{18}{2}{1:49}. Forward-looking financial markets do the same: changed expectations change asset prices (standard finance), and the
money view adds that \emph{price changes cause cash flows today}, tightening or loosening the survival constraint \ts{18}{2}{2:30}.
The dramatic case: the second round of funding is refused, or the bond meant to take out short-term funding cannot be sold, ``then I'm dead''; the everyday case:
pressure on money market prices \ts{18}{2}{4:39}.

% ------------------------------------------------------------------
\section{FT: Argentina in court}\label{sec:ff-ft}

\begin{casestudy}[FT, November 2012: ``Argentina in court to fight debt ruling'']
In 2001 Argentina defaulted. Eventually 93\% of bondholders by value accepted a restructuring (a reduction of the debt) and got new bonds; the other 7\%,
holdouts, mostly hedge funds, sued for full payment \ts{18}{1}{0:07}. A New York court ruled that unless the
holdouts are paid in full, Argentina may not pay the restructured bonds, since all bondholders must be treated equally; this undermines restructuring itself
\ts{18}{1}{1:29}. With a payment on the restructured bonds due the following month, Argentina faces a ``monetary cliff'': pay the
holdouts too, or default on everything \ts{18}{1}{2:51}. The FT's editorial underlines the importance of collective action clauses: one holdout
blocking agreement is a social loss, and taking a loss can be in the bondholders' collective interest \ts{18}{1}{3:12}.

The article in the course file (Jude Webber, \emph{FT}, 26 November 2012, p.~6; \file{materials/readings/L18-Financial_Times_-_Presscuttings.pdf}) adds the details:
\begin{itemize}
  \item The 2001 default was on nearly \$100bn. The restructurings were in 2005 and 2010.
  \item Judge Thomas Griesa ordered Argentina to pay \$1.3bn to the holdouts when it paid the restructured bonds on 15 December, and to deposit the money in escrow pending appeal. Third parties helping it evade the order would count as aiders and abettors.
  \item The dispute turns on the ``pari passu'' (``equal step'') clause.
  \item Argentina's own law forbade reopening the restructuring. Failing to comply could mean a technical default ``perhaps as early as mid-January''.
\end{itemize}
\end{casestudy}

\begin{definition}[New concepts: restructuring, holdout, collective action clause]
\emph{Restructuring}: renegotiating a debt to less (longer terms, lower rates, lower principal). \emph{Holdouts}: creditors who refuse the deal. A \emph{collective
action clause} (CAC) lets a qualified majority of bondholders bind the minority, so there can be no holdouts; the 2001 Argentine bonds, issued under New York law,
had none \ts{18}{1}{0:48}.
\end{definition}

\begin{history}[The end of the Argentine holdout saga]
The ruling was by Judge Thomas Griesa (S.D.N.Y.), interpreting the bonds' \emph{pari passu} clause; it was upheld on appeal, and in June 2014 the US Supreme Court
declined to hear Argentina's appeal. Argentina defaulted again in July 2014 and settled with most holdouts in 2016.\par
\emph{References:} \emph{NML Capital, Ltd.\ v.\ Republic of Argentina}, 699 F.3d 246 (2d Cir.\ 2012); IMF, ``Strengthening the Contractual Framework to Address
Collective Action Problems in Sovereign Debt Restructuring'', 2014.
\end{history}

% ------------------------------------------------------------------
\section{Forwards versus futures}\label{sec:ff-fwdfut}

\paragraph{A forward is advance clearing.} Recall the forward rate agreement as a swap of IOUs between a bank and firm A: a six-month loan against a three-month
deposit, notional amounts, no money changing hands, locking in a three-month rate three months ahead, $f_{3,6}$, defined by forward interest parity
\ts{18}{3}{0:24}:
\begin{keyformulas}
\textbf{Forward interest parity} \ts{18}{3}{1:50}: $(1+r_{0,3})(1+f_{3,6}) = 1+r_{0,6}$, where $r_{0,3}$ and $r_{0,6}$ are today's three- and six-month
rates: the forward rate makes the long rate equal to rolling over two short ones.
\end{keyformulas}
The bank, having promised to lend in three months, hedges with a firm B wanting the opposite, to lock in a deposit rate: the bank's positions offset. In practice
these are forward contracts listed off balance sheet, not parallel loans \ts{18}{3}{2:54}. The
bank has \emph{cleared in advance}: in three months A will be a deficit agent and B a surplus agent, and the flows are already agreed; one way to evade the survival
constraint is to take care of it ahead of time \ts{18}{3}{4:40}. The bank runs a matched book and earns the bid--ask spread
\ts{18}{3}{5:42}.

\paragraph{From matched book to speculators.} There is no reason for equal numbers of fundamental customers on both sides; imbalance pushes prices until dealers
take the other side onto their balance sheets. When the banking system as a whole has as much risk as it wants, banks find hedge funds and speculators in the
\emph{futures} market: a speculator goes long futures to help the bank square its book with a short futures position
\ts{18}{3}{6:22}. The chain: client $\to$ bank $\to$ other banks (FRAs) $\to$ futures exchange
(speculators) $\to$ spot money market, where a rate is realized in three months \ts{18}{3}{8:07}.

\begin{nfigure}
\begin{tikzpicture}[font=\scriptsize,node distance=1.2cm,box/.style={draw,rounded corners,align=center,minimum width=2.1cm,minimum height=0.8cm}]
  \node[box] (c) {client (firm A)\\locks in a loan};
  \node[box,right=0.6cm of c] (b) {bank\\(matched book?)};
  \node[box,right=0.6cm of b] (o) {other banks\\(FRAs)};
  \node[box,right=0.6cm of o] (f) {futures exchange\\(speculators)};
  \node[box,right=0.6cm of f] (s) {spot money\\market at $t=3$};
  \draw[-{Stealth}] (c) -- (b); \draw[-{Stealth}] (b) -- (o); \draw[-{Stealth}] (o) -- (f); \draw[-{Stealth}] (f) -- (s);
  \node[below=0.35cm of o,align=center] {$f_{3,6} \;>\; \text{futures rate} \;>\; \E[r_{3,6}]$: each hedge takes a piece of the premium};
\end{tikzpicture}
\caption{A sequence of hedging transactions \ts{18}{3}{8:31}.}\label{fig:ff-chain}
\end{nfigure}

\paragraph{A stylized fact.} The forward rate is not an unbiased forecast of the future spot rate (the expectations hypothesis fails), and one reason, in this
course, is a mismatched book in the forward market: dealers must be paid to take the risk \ts{18}{3}{9:32}. The sequence of hedges
makes sense only if the forward rate exceeds the expected spot rate, each taking a piece; so one also expects a gap between forward and futures rates
\ts{18}{3}{10:33}:
\begin{keyformulas}
\textbf{Stylized fact (as stated in class)} \ts{18}{3}{11:14}: \quad $f_{3,6} \;>\; \text{futures rate} \;>\; \E[r_{3,6}]$. The forward--futures gap
is small compared with the forward--expected spot gap; finance usually attributes it to something other than liquidity.
\end{keyformulas}

\begin{reading}[Hicks, \emph{Value and Capital} (1939), ch.~XI ``Interest'', pp.~141--152]
The course reading (\file{materials/readings/L18-Hicks_Value_and_Capital.pdf}) is the classic source of this stylized fact.
\begin{itemize}
  \item A loan is a transaction with only one side executed now. Any loan can be reduced to a money loan plus spot and forward transactions. With forward markets, ``own'' rates of interest are implicitly fixed: with money at 5\% and coffee futures 3\% above spot, the coffee rate is about 2\%.
  \item Long rates are built up from \emph{forward} short rates, determined like commodity futures prices by hedgers and speculators.
  \item The forward market for loans has a ``constitutional weakness'' on one side: most people prefer to lend short. Only speculators fill the gap, so the forward rate must exceed the rate they expect, by a risk premium which ``corresponds exactly to the `normal backwardation' of the commodity markets'' (p.~147).
  \item Hence, if short rates are not expected to change, forward rates exceed current short rates by this premium.
\end{itemize}
This is the course's mismatched forward book, with the dealer's required compensation as Hicks's premium.
\end{reading}

\paragraph{The difference is cash flow.} A forward causes no cash flow until maturity; futures are \emph{marked to market} and cause cash flows every day, in or
out, as prices move. So do most derivatives, through margin and collateral calls, ``and those collateral calls cause people to go out of business'': the survival
constraint \ts{18}{3}{12:16}.

\begin{definition}[New concepts: futures contract, marking to market]
A \emph{futures contract} is a standardized forward traded on an exchange, whose price is reset every day; \emph{marking to market} means settling the change in
value in cash each day (through margin), so that the contract's value returns to zero \ts{18}{3}{12:38}\ts{18}{5}{0:04}. (\emph{Futures}: short for contracts for future delivery.)
\end{definition}

% ------------------------------------------------------------------
\section{The value of a forward over its life}\label{sec:ff-value}

\paragraph{Convergence.} Rewrite forward interest parity as \ts{18}{4}{0:03}
\[
  \frac{1}{1+f^{0}_{3,6}} = \frac{1+r_{0,3}}{1+r_{0,6}}, \qquad \frac{1}{1+f^{1}_{3,6}} = \frac{1+r_{1,3}}{1+r_{1,6}},\ \ldots,\ \frac{1}{1+f^{3}_{3,6}} = \frac{1}{1+r_{3,6}},
\]
where $f^{t}_{3,6}$ is the forward rate quoted at time $t$; in general they all differ. At maturity the forward rate equals the spot rate
\ts{18}{4}{1:11}. So the contract between firm A and the bank, a zero-value swap of
IOUs at inception, changes value: a zero-sum game in which, if rates behave as usual, the bank wins, having locked in a lending rate above the expected spot rate;
the firm pays a premium for peace of mind about a serious survival constraint \ts{18}{4}{3:22}.

\begin{intuition}[Coordinating views of the future]
In the notes' framing \mn{18}{1}: in finance the future determines the present, but nobody knows the future, so there are many views of it. How is one path chosen and diverse views
coordinated? By market pricing of the views, and by the effect of prices on behaviour through the survival constraint. At one extreme, if the market changes its mind about your view you
cannot roll your funding; more subtly, changing ideas about the future change cash flows today, loosening the constraint for some and tightening it for others, as futures prices do.
\end{intuition}

The classic forward is the wheat farmer, naturally long wheat, and the baker, naturally short: by a forward contract both lock in the price, each hedging a natural exposure; at delivery one
side ``wins'', but in a forward the winnings are only notional until the end \mn{18}{3}. The forward price is an arbitrage condition: if $K>S_0e^{rT}$, buy the bond spot and sell it forward,
locking in a return above the borrowing rate $r$; if $K<S_0e^{rT}$, sell spot and buy forward \mn{18}{4}. (The notes follow J.~Hull, \emph{Options, Futures, and Other Derivatives}, 5th ed.,
equations 3.5 and 3.9.)

\paragraph{Continuous-time version.} In price terms, with $K$ the delivery (forward) price, $S_t$ the spot price of the underlying (here the six-month bill) and $r$
the interest rate \ts{18}{4}{5:46}:
\begin{keyformulas}
\textbf{Forward price and value} \ts{18}{4}{7:10}:
\[
  K = S_0\,e^{rT} \quad\text{(1)}, \qquad f_t = S_t - K\,e^{-r(T-t)} \quad\text{(2)}, \qquad F_t = S_t\,e^{r(T-t)} \quad\text{(3)}.
\]
$K$ is fixed at inception; the value $f_t$ of the forward to the long side moves with the spot price and with the shrinking discount as maturity approaches:
$f_0=0$ and $f_T = S_T - K$. The futures price $F_t$ (3) is $K$ reset every day so that the contract's value is again zero \ts{18}{5}{0:04}.
\end{keyformulas}

Banking deals constantly with zero-value contracts that later acquire value, ``a little mind-blowing'' \ts{18}{4}{10:00}. The
same apparatus applies to commodities with storage costs and to coupon bonds with more complicated formulas; here we keep discount bonds of the money market
\ts{18}{4}{11:28}. With a forward there are no cash flows until the end, when the shorts pay the longs $S_T-K$ (in financial
contracts by cash netting) or the reverse \ts{18}{4}{12:30}.

\begin{nfigure}
\begin{tikzpicture}[font=\scriptsize,x=1cm,y=1cm]
  \draw[-{Stealth[length=4pt]}] (0,0) -- (8.3,0) node[right]{$t$};
  \draw[-{Stealth[length=4pt]}] (0,-1.6) -- (0,2.2) node[above]{$f_t$ (value to the long side)};
  \draw[dashed] (7.5,-1.6) -- (7.5,2.0) node[above]{$T$};
  \draw[very thick,defcol] (0,0) .. controls (1,1.2) and (1.8,1.4) .. (2.5,0.6) .. controls (3.3,-0.4) and (4,-1.3) .. (4.8,-0.6) .. controls (5.6,0.2) and (6.6,1.2) .. (7.5,1.0);
  \fill (7.5,1.0) circle (1.5pt) node[right,align=left]{$S_T-K$: the only\\cash flow of a forward};
  \node[align=left] at (1.5,1.75) {futures: cash to longs};
  \node[align=left] at (4.0,-1.45) {futures: cash to shorts};
  \node[align=left] at (6.4,1.7) {to longs};
\end{tikzpicture}
\caption{The same price history, two cash-flow profiles: a forward pays once at maturity, a futures contract pays every day as the value changes
\ts{18}{4}{11:01}\ts{18}{5}{1:47}. Schematic.}\label{fig:ff-paths}
\end{nfigure}

\paragraph{Margin.} The payments go through margin accounts at the futures clearinghouse, and both sides post margin, since at inception both are equally likely to lose; the clearinghouse
accepts liquid securities as margin, repriced every day, so collateral and contract are both marked to market. The cumulative futures payments equal the forward's final payment, but they come
every day: ``a very concrete way in which views about the future are settled today'' \mn{18}{6}. (The notes' formula abstracts from interest on the fluctuating margin balances \mn{18}{5}.)

\paragraph{Futures cash flows.} Whenever the futures price changes, cash moves from losers to winners: while the forward's value rises, from the shorts to the
longs; when it falls, back \ts{18}{5}{1:06}. With a forward, who was winning in the middle does
not matter; with futures the largest swings may come in the middle, and you must be ready to absorb them. Futures carry liquidity risk that forwards do not, which
may be why their prices differ, though the formulas treat them as interchangeable \ts{18}{5}{2:50}.

% ------------------------------------------------------------------
\section{Cash and carry arbitrage}\label{sec:ff-cashcarry}

\begin{aside}[Ten years of worrying]
Stigum's pages 718--722 show real cash and carry trades making money, which, by finance theory, should not be possible. ``This lecture is a consequence of 10 years
of me figuring this stuff out'': find something that does not seem right and keep worrying at it \ts{18}{5}{4:12}.
\end{aside}

\paragraph{The trade.} Buy a six-month Treasury bill and finance it with a three-month repo: borrowing short, lending long; the price risk (the value of the then
three-month bill in three months) is hedged by a short futures position \ts{18}{6}{0:03}.
Long bill plus short repo is a synthetic forward (the six-month loan against the three-month deposit of the balance sheets): so the position is \emph{long forward,
short futures} for the same term. Why should it make any profit? \ts{18}{6}{2:08}

\begin{nfigure}
\begin{taccts}
\tacct[2.6cm]{Cash and carry position}{6-month Treasury bill}{3-month repo\\Short futures}
\end{taccts}
\caption{Stigum's cash and carry: long positions above, short positions below the line in class; long bill $+$ short repo $=$ synthetic long forward
\ts{18}{6}{1:48}.}\label{fig:ff-cc}
\end{nfigure}

\begin{keyformulas}
\textbf{Full carry pricing and the implied repo rate} \ts{18}{6}{3:51}:
\[
  F_0 = S_0\,e^{rT}\ \text{(full carry)}, \quad F_0 > S_0\,e^{rT}\ \text{(cash and carry)}, \quad F_0 < S_0\,e^{rT}\ \text{(reverse).}
\]
The \emph{implied repo rate} $\rho$ is defined by $F_0 = S_0\,e^{\rho T}$: the rate that would make full carry pricing true. Cash and carry ($\rho>r$) is like
borrowing at the actual repo rate $r$ and lending at $\rho$: an interest rate spread.
\end{keyformulas}

\begin{reading}[Stigum's example, \emph{The Money Market}, pp.~718--722]
Stigum's own numbers (\file{materials/readings/L18-Pages_from_Stigum.pdf}) are dated 28 October 1982, ``when liquidity in T-bills was at its peak''.
\begin{enumerate}
  \item Buy \$1 million of the 24 March 1983 bill, 147 days to maturity, at a discount rate of 8.26: price \$966{,}271.67.
  \item Sell the December 1982 bill future, which expires in 56 days, at 91.72 (rate 8.28). Delivering the bill, then a 91-day bill, brings in \$979{,}070.00.
  \item The holding-period yield over the 56 days, on a 360-day basis, is the \emph{implied repo rate}: 8.51\%.
  \item With 56-day term repo at 8.25, the trade earns 26\,bp. On \$10 million that is \$3{,}908.03.
\end{enumerate}
Stigum lists the ``slips twixt the cup and the lip'': commission, margin calls, convergence-price risk and back-office costs. He judges margin calls a small threat. Even a 100\,bp rally on the first day would cost only 2\textonequarter\,bp of the spread. In Mehrling's reading this cash-flow timing is the whole point: what is left after the price risk is hedged is liquidity risk.

(Table 15.4 dates the futures sale ``March 28, 1982''; the context shows it is 28 October 1982.)
\end{reading}

\paragraph{Explained as liquidity risk.} Look at the cash flows \ts{18}{7}{0:03}:
\begin{itemize}
  \item \textbf{long forward}: since the expectations hypothesis fails, the long side probably wins, but receives cash only at $T$;
  \item \textbf{short futures}: probably loses, and pays all along, with margin calls until $T$; if you cannot keep up, the clearinghouse liquidates your
        position \ts{18}{7}{1:45}.
\end{itemize}
All price risk is gone; what is left is liquidity risk: the inflows will cover the outflows, ``but not at the same time''. The position is risk-free only if one could
somehow use the certain final inflow to pay the earlier outflows \ts{18}{7}{2:26}. The same banking view that explains the failure
of EH (an imbalance in the forward market) explains Stigum's confusing pages: forwards and futures do not trade at the same price \ts{18}{7}{3:28}. The notes stress that the negative
flows might all come at once: the volatility of the bill's spot price is what creates liquidity risk, and the two anomalies, the failure of EH and the profit of cash and carry, have the same origin,
a mismatch in the natural forward interest rate positions of the real economy. A former teaching assistant of the course, Daniel Neilson, showed how forward--futures deviations can serve as
an observable proxy for the unobservable deviation of forward rates from expected spot rates, a test of the liquidity premium theory of the term structure \mn{18}{8}.

\begin{coursenote}
Mehrling notes that the standard explanation, reinvestment of daily futures cash flows, ``goes in the wrong direction'' and may even hide the size of the liquidity
premium \ts{18}{7}{4:12}. In the literature this is the forward--futures price difference under stochastic interest
rates (J.~C.~Cox, J.~E.~Ingersoll and S.~A.~Ross, ``The relation between forward prices and futures prices'', \emph{Journal of Financial Economics} 9, 1981); for
interest-rate futures it implies a futures rate \emph{above} the forward rate (the ``convexity adjustment''), the opposite of the class's stylized fact
$f>$ futures.
\end{coursenote}

\paragraph{Or counterparty risk?} A student suggests credit risk: a forward with a repo counterparty is one big payment at the end, never sure until paid, while
futures are margined by an exchange, with warning from many small payments. More credit risk in the forward would call for compensation in the right direction
\ts{18}{8}{0:03}. Typical of these debates: a spread exists, and one must list all candidate causes
(cash-flow timing, counterparty risk, reinvestment) and do the empirical work to separate them; probably traders have done it and are making money, but ``it's not
been published in the academic literature'' \ts{18}{8}{1:46}.

\paragraph{A natural banking business.} Another student: futures need cash at a moment's notice, so how does that affect banks? Banks use futures to hedge forward
positions, so value fluctuations offset; the naked speculator on the other side must post margin. What the bank is not hedged against is the liquidity risk, but a
bank has the discount window and interbank relationships, a comparative advantage in liquidity risk, so cash and carry is natural banking business; speculators bear
liquidity risk and are paid the term premium, the difference between the forward rate and the expected spot rate
\ts{18}{9}{0:01}.

% ------------------------------------------------------------------
\flushside
\section*{Key formulas and ideas}
\begin{keyformulas}
\begin{tabularx}{\linewidth}{@{}l L@{}}
Advance clearing & forward contracts settle future payments today; asset price changes cause cash flows today\\
Forward parity & $(1+r_{0,3})(1+f_{3,6})=1+r_{0,6}$; at maturity the forward rate equals the spot rate\\
Forward & $K=S_0e^{rT}$, $f_t=S_t-Ke^{-r(T-t)}$; one cash flow $S_T-K$ at $T$\\
Futures & $F_t=S_te^{r(T-t)}$ reset daily; marked to market, cash flows all along\\
Stylized fact & $f_{3,6}>$ futures rate $>\E[r_{3,6}]$: each hedger takes part of the liquidity premium\\
Cash and carry & long bill $+$ short repo $+$ short futures; profit if $\rho>r$; the residual risk is liquidity risk\\
\end{tabularx}
\end{keyformulas}

\section*{Exercises}

\begin{exercise}[Forward interest parity]
Three-month and six-month rates are 1.0\% and 2.2\% (not annualized). Compute $f_{3,6}$. One month later the two-month and five-month rates are 0.7\% and 1.8\%: compute
the new forward rate and say whether the bank that lent forward has gained or lost.
\end{exercise}

\begin{exercise}[Forward versus futures cash flows]
A forward and a futures contract on the same bill have $K=F_0=98$. The futures price then goes 98.5, 97.0, 96.5, 98.0, 99.0 at maturity. Compute the daily cash
flows of a short futures position and the single cash flow of a short forward. What is the largest cumulative cash outflow of the futures position?
\end{exercise}

\begin{exercise}[Implied repo rate]
A six-month bill trades at $S_0=99$ and the three-month futures price on it is $F_0=99.6$; the three-month repo rate is 2\% a year (continuous compounding). Compute the
implied repo rate. Is a cash and carry trade profitable, and what liquidity must the trader hold?
\end{exercise}

\begin{exercise}[Who should do it?]
Explain why a bank with access to the discount window, rather than a hedge fund, is the natural holder of a cash and carry position, and what happens to the
forward--futures spread if banks' access to liquidity is restricted.
\end{exercise}
